\section{From a small matrix projection to a scalar projection}
\label{sec:scalar}

We first isolate a reduction that will be useful independently of the
construction. Its trace convention is essential.
For related results on the scale and comparison of projections in
reduced free products, see Dykema--R{\o}rdam \cite{DR1998,DR2000,DRII2000}.
The explicit row estimate used here is proved below.
For a discrete group \(\Lambda\), let
\[
 (\lambda_\Lambda(g)\xi)(h)=\xi(g^{-1}h),\qquad
 C_r^*(\Lambda)=
 \overline{\lambda_\Lambda(\mathbb C[\Lambda])}^{\,\|\cdot\|}.
\]
The canonical trace is
\(\tau_\Lambda(x)=\langle x\delta_1,\delta_1\rangle\).
It is tracial on finite group-algebra sums by taking the identity
coefficient, and hence on their norm closure. It is faithful:
if \(x\ge0\) and \(\tau_\Lambda(x)=0\), then
\(x^{1/2}\delta_1=0\). Since \(x^{1/2}\) commutes with the
right regular action, it annihilates every basis vector and is zero.
On matrices we always use the \emph{unnormalized} extension
\begin{equation}
 \tau_{\Lambda,N}(x)=\sum_{j=1}^N\tau_\Lambda(x_{jj}).
 \label{scalar:matrix-trace}
\end{equation}
The same formula applies to positive equivariant operators on finite
regular multiples. In particular,
\(\tau_{\Lambda,N}(p)>0\) for a nonzero projection \(p\).

\begin{lemma}[Reduced subgroup inclusion]
\label{scalar:subgroup}
If \(\Lambda\) is a subgroup of a discrete group \(K\), the inclusion of
group algebras extends to an isometric unital inclusion
\(C_r^*(\Lambda)\subseteq C_r^*(K)\), preserving the canonical traces
and their finite matrix extensions.
\end{lemma}

\begin{proof}
Choose one representative \(v\) from every left \(\Lambda\)-orbit
\(\Lambda v\) in \(K\). The unitary
\(\delta_\gamma\mapsto\delta_{\gamma v}\) identifies the restriction
of the left regular action to \(\ell^2(\Lambda v)\) with the left
regular action on \(\ell^2(\Lambda)\). Thus the representation
restricted to \(\mathbb C[\Lambda]\) is a direct sum of regular
representations, so its norm is precisely the reduced norm on
\(\Lambda\). Taking closures gives the inclusion. The identity
diagonal proves preservation of trace; summing diagonal entries proves
the matrix statement.
\end{proof}

\begin{proposition}[Scalar compression after a free product]
\label{scalar:compression}
Let \(\Lambda\) be a discrete group and
\(p\in M_N(C_r^*(\Lambda))\) a projection with
\[
 0<a:=\tau_{\Lambda,N}(p)<\frac12.
\]
Put \(K=\Lambda*\mathbb F_N\), where \(t_1,\ldots,t_N\) freely
generate \(\mathbb F_N\), and regard \(p\) as a matrix over
\(C_r^*(K)\) by Lemma~\ref{scalar:subgroup}. Set
\[
 U=(\lambda_K(t_1),\ldots,\lambda_K(t_N)),\qquad
 A=pU^*Up.
\]
Then \(A\) is invertible in the unital corner
\(pM_N(C_r^*(K))p\), and
\begin{equation}
 e=(Up)A^{-1}(Up)^*
 \label{scalar:projection-formula}
\end{equation}
is a projection in the scalar algebra \(C_r^*(K)\) with
\(\tau_K(e)=a\).
If \(\Lambda\) is finitely generated and torsion-free, then so is \(K\).
\end{proposition}

\begin{proof}
To bound \(U\) below on \(p\ell^2(K)^N\), we separate
noncancelling inputs, whose outputs are disjoint, from cancelling
inputs, whose size both before and after applying \(U\) is controlled
by the trace \(a\).
Use free-product normal forms, writing entries of the free factor
in reduced free letters. Let \(\mathcal R\) consist of the identity
and the normal words whose first factor is in \(\mathbb F_N\).
Every element of \(K\) has a unique expression \(\gamma v\) with
\(\gamma\in\Lambda\) and \(v\in\mathcal R\): remove its initial
\(\Lambda\)-factor, if present. Thus \(\mathcal R\) is a transversal
for the orbits used in Lemma~\ref{scalar:subgroup}.

Write \(a_j=\tau_\Lambda(p_{jj})\), so
\(a_j\ge0\) and \(\sum_j a_j=a\). If
\(\xi=p\xi\in\ell^2(K)^N\), put
\[
 m_v=\|\xi|_{\Lambda v}\|^2,\qquad
 \sum_{v\in\mathcal R}m_v=\|\xi\|^2.
\]
Restriction to an orbit commutes with \(p\).
Applying Cauchy--Schwarz to the projection on that orbit gives
\begin{equation}
 |\xi_j(v)|^2
 =|\langle \xi|_{\Lambda v},p\delta_{(v,j)}\rangle|^2
 \le a_jm_v.
 \label{scalar:pointwise}
\end{equation}
Here
\(\|p\delta_{(v,j)}\|^2
=\langle p\delta_{(v,j)},\delta_{(v,j)}\rangle=a_j\).

In coordinate \(j\), call a position bad if its reduced word starts
with \(t_j^{-1}\), and good otherwise. Split \(\xi\) accordingly as
\(\xi=\xi_{\mathrm g}+\xi_{\mathrm b}\).
Every bad position belongs to \(\mathcal R\); its first letter
determines its bad coordinate uniquely. Therefore
\begin{equation}
 \|\xi_{\mathrm b}\|^2
 \le a\sum_{v\in\mathcal R}m_v=a\|\xi\|^2.
 \label{scalar:bad-input}
\end{equation}
Left multiplication by \(t_j\) sends every good position to a word
starting with \(t_j\). These output sets are disjoint for different
\(j\), and multiplication is injective within each coordinate. Hence
\begin{equation}
 \|U\xi_{\mathrm g}\|^2
 =\|\xi_{\mathrm g}\|^2
 \ge(1-a)\|\xi\|^2.
 \label{scalar:good-output}
\end{equation}

Bad outputs may coincide, and require a separate estimate.
For an output word \(h\), let \(J(h)\) be the coordinates whose bad
input \(v_{h,j}=t_j^{-1}h\) contributes to \(h\).
By \eqref{scalar:pointwise},
\[
 |(U\xi_{\mathrm b})(h)|^2
 \le\left(\sum_{j\in J(h)}a_j\right)
      \left(\sum_{j\in J(h)}m_{v_{h,j}}\right)
 \le a\sum_{j\in J(h)}m_{v_{h,j}}.
\]
The map \((h,j)\mapsto v_{h,j}\), restricted to these bad inputs,
is injective: the first letter of \(v_{h,j}\) determines \(j\),
and then \(h=t_jv_{h,j}\). Summing over \(h\) gives
\begin{equation}
 \|U\xi_{\mathrm b}\|^2\le a\|\xi\|^2.
 \label{scalar:bad-output}
\end{equation}
Equations~\eqref{scalar:good-output} and
\eqref{scalar:bad-output}, and the triangle inequality, imply
\begin{equation}
 \|U\xi\|\ge c_a\|\xi\|,\qquad
 c_a=\sqrt{1-a}-\sqrt a>0.
 \label{scalar:lower-bound}
\end{equation}

Consequently \(A\ge c_a^2p\) in its C*-corner.
Also \(UU^*=N1\), so \(A\le Np\). In particular, its inverse belongs
to the corner itself: explicitly,
\[
 A^{-1}=\frac1N\sum_{k=0}^\infty(p-A/N)^k
\]
converges in norm because
\(\|p-A/N\|\le1-c_a^2/N<1\).
Continuous functional calculus also supplies \(A^{-1/2}\) there.
The row \(V=UpA^{-1/2}\) belongs to
\(M_{1,N}(C_r^*(K))\) and satisfies \(V^*V=p\).
Thus \(e=VV^*\) is the scalar projection in
\eqref{scalar:projection-formula}. Writing \(V=(V_1,\ldots,V_N)\)
and using traciality,
\[
 \tau_K(e)
 =\sum_j\tau_K(V_jV_j^*)
 =\sum_j\tau_K(V_j^*V_j)
 =\tau_{K,N}(p)=a.
\]
It is therefore neither zero nor one.

Finally, adjoining the finite set \(t_1,\ldots,t_N\) preserves
finite generation. A nontrivial element of a free product is
conjugate either to a nontrivial factor element or to a cyclically
reduced word with at least two alternating factor entries.
Positive powers of the latter have strictly increasing reduced
length. If both factors are torsion-free, neither case has finite
order. This proves the last assertion.
\end{proof}

It remains to construct a finitely generated torsion-free
\(\Gbase\) and a projection satisfying the small
\emph{unnormalized} trace hypothesis of
Proposition~\ref{scalar:compression}. The next sections do so.
